%% =====================================================================
%%  Avant-propos
%% =====================================================================
\thispagestyle{otFront}
\phantomsection\pdfbookmark[0]{Avant-propos}{bm-avant-propos}
\otPageTitre{Avant-propos}{}

\begin{minipage}{128mm}
{\fontsize{14}{20}\selectfont\itshape\color{otNavy}
Ce livre a été conçu pour vous accompagner tout au long de votre première
année de classe préparatoire MPSI, du premier cours aux révisions d'été.\par}
\vspace{7mm}

\fontsize{11}{16}\selectfont
Il couvre l'intégralité du programme officiel de mathématiques des classes
préparatoires MPSI et MP2I (BO spécial n\textsuperscript{o}\,1 du 11 février
2021), dans l'ordre des deux semestres : vingt chapitres, du raisonnement et
du vocabulaire ensembliste aux fonctions de deux variables.\par\vspace{3.5mm}

Chaque chapitre s'ouvre sur \textbf{L'essentiel} : une synthèse structurée du
cours — définitions, théorèmes, démonstrations à connaître, méthodes et
pièges à éviter — pour réviser efficacement avant une colle ou un devoir
surveillé. Elle ne remplace pas le cours de votre professeur : elle vous aide
à le retenir et à le mobiliser.\par\vspace{3.5mm}

Viennent ensuite vingt-cinq exercices par chapitre, classés en quatre
rubriques. Les \textbf{Applications} vérifient que les définitions et les
théorèmes sont en place ; l'\textbf{Entraînement} installe les méthodes
incontournables, celles que l'on retrouve en colle et en devoir surveillé ;
l'\textbf{Approfondissement} propose des problèmes plus ambitieux, dans
l'esprit des écrits de concours ; le \textbf{Défi}, enfin, conduit pas à pas
jusqu'à un résultat célèbre.\par\vspace{3.5mm}

La partie \textbf{Évaluations} rassemble seize sujets inédits : devoirs
surveillés, devoirs maison, examens et sujets type concours — CCINP,
Mines-Ponts, Centrale-Supélec, X-ENS. Tous les exercices et tous les sujets
sont corrigés en détail en fin d'ouvrage : les corrigés indiquent, à chaque
étape, la méthode utilisée et les réflexes de rédaction attendus en prépa.
Cherchez toujours par vous-même avant de les consulter : c'est en cherchant
que l'on progresse.\par\vspace{3.5mm}

Le tableau de bord, au début du livre, vous permet de suivre votre
progression chapitre par chapitre et de repérer, avant chaque colle, les
exercices à retravailler.\par\vspace{9mm}

{\otMedium\fontsize{10}{13}\selectfont Bon travail, et bonne réussite en MPSI !}\par
\vspace{3mm}
{\otLight\fontsize{9}{12}\selectfont\color{otMuted}L'équipe pédagogique OptiTAB}\par
\vspace{8mm}
\includegraphics[height=9mm]{icone-couleur}
\end{minipage}
\clearpage
